\subsection{Definicões}
\begin{definition}[ponto interior]
  Um $a\in X\subseteq \mathbb{R}^{n}$ é interior a $X$ se existe $r>0$ tal que $B(a;r)\subseteq X$. 
\end{definition}
\begin{definition}[Interior do $A$]
  Dado um conjunto $A\subseteq \mathbb{R}^{n}$ definimos
  \begin{align*}
    int(A):=\left\{x\in A: x\text{ e um ponto interior do }A\right\}
  \end{align*}
\end{definition}
\begin{definition}[Bola aberta]
  Seja $a\in \mathbb{R}^{n}$ e $r>0$, então 
  \begin{align*}
    B(a;r):=\left\{x\in \mathbb{R}^{n}:|x-a|<r\right\}
  \end{align*}
\end{definition}
\begin{definition}[Bola fechada]
  Seja $a\in \mathbb{R}^{n}$ e $r>0$, então 
  \begin{align*}
    B(a;r):=\left\{x\in \mathbb{R}^{n}:|x-a|\leq r\right\}
  \end{align*}
\end{definition}
\begin{definition}[Esfera]
  Seja $a\in \mathbb{R}^{n}$ e $r>0$, então 
  \begin{align*}
    B(a;r):=\left\{x\in \mathbb{R}^{n}:|x-a|=r\right\}
  \end{align*}
\end{definition}
\begin{example}[Norma Euclidiana]
  Definimos a norma Euclidiana como $|x|=\sqrt{\sum_{i=1}^{n}x_i^2}$, o gráfico da bola é como
  \begin{center}
    \begin{tikzpicture}[scale=0.8]

  % Ejes
  \draw[->] (-3,0) -- (3,0) node[right] {$x_1$};
  \draw[->] (0,-3) -- (0,3) node[above] {$x_2$};
  
  % Bola de radio r
  \draw[dotted] (0,0) circle (2);
      
  % Radio
  \draw[->] (0,0) -- (2,0);
  \node[below] at (1,0) {$r$};
     
  % Centro
  \fill (0,0) circle (2pt);
  \node[below left] at (0,0) {$0$};
\end{tikzpicture} 
     
  \end{center}
\end{example}
\begin{example}[Norma da Suma]
  Definimos a norma da suma como $|x|_{S}=\sum_{i=1}^{n}x_i$, o gráfico da bola é como
  \begin{center}
    \begin{tikzpicture}[scale=0.8]

  % Ejes
  \draw[->] (-3,0) -- (3,0) node[right] {$x_1$};
  \draw[->] (0,-3) -- (0,3) node[above] {$x_2$};
      
  % Bola L1
  \draw[dotted] (2,0) -- (0,2) -- (-2,0) -- (0,-2) -- cycle;
      
  % Radio
  \draw[->] (0,0) -- (2,0);
  \node[below] at (1,0) {$r$};

  % Centro
  \fill (0,0) circle (2pt);
\end{tikzpicture}  


  \end{center}
\end{example}
\begin{example}[Norma da Máximo]
  Definimos a norma da máximo como $\displaystyle |x|_{M}=\max_{i=1,\cdots,n}|x_i|$, o gráfico da bola é como
  \begin{center}
    \input{images/bolamaximo.tex}
  \end{center}
\end{example}
\begin{note}
  As normas $|\phantom{x}|,|\phantom{x}|_{S},|\phantom{x}|_{M}$ são equivalenten no seguinte sentido
  \begin{align*}
    |x|_{M}\leq|x|\leq|x|_{S}\leq n|x|_{M}
  \end{align*}
\end{note}
\begin{definition}[Conjunto limitado]
  Um conjunto $X\in\mathbb{R}^{n}$ é limitado se existe uma bola que ó contém. 
\end{definition}
\begin{definition}[Conjunto aberto]
  Um conjunto $A$ é aberto se todos os seus pontos são interiores, i,e. $A=int(A)$. 
\end{definition}
\subsection{Propriedades dos abertos}
\begin{theorem}
  A interseção finita de abertos é aberta.
\end{theorem}
\begin{proof} 
  Suponha $X_i\subseteq \mathbb{R}^{n}$ com $i=1,\cdots,k$ conjuntos abertos, então dado $x\in \bigcap_{i=1}^{k}X_i$ sabemos que $x\in X_i$ para tudo $i=1,\cdots,=k$, logo como $x$ é ponto interior de cada conjunto existem $r_i>0$ de modo que $x\in B(x,r_i)\subseteq X_i$.\\
  Suponha $\displaystyle \delta=\min_{i=1,\cdots,k} r_i$, assim por propriedades do mínimo sabemos que $B(x,\delta)\subseteq X_i$ para tudo $i=1,\cdots,k$, portanto $x\in B(x,\delta)\subseteq \bigcap_{i=1}^{k}X_i$, i,e, $x\in Int\left( \bigcap_{i=1}^{k}X_i \right)$ e como $x$ é arbitrario podemos concluir que $\bigcap_{i=1}^{k}X_i=int\left( \bigcap_{i=1}^{k}X_i \right)$.
\end{proof}
\begin{theorem}
  A reunión arbitraria de abertos é aberta.
\end{theorem}
\begin{proof} 
  Suponha $X_i\subseteq \mathbb{R}^{n}$ com $i\in \Lambda$ com $\Lambda$ índice arbitrario, então dado $x\in \bigcup_{i\in\Lambda}X_i$, sabemos que existe pelo menos um $j\in\Lambda$ tal que $x\in X_j$, logo $x$ é ponto interior de $X_j$ e portanto existe $r>0$ tal que $x\in B(x,r)\subseteq X_j \subseteq \bigcup_{i\in\Lambda}X_i$ e portanto $x\in int\left( \bigcup_{i\in\Lambda}X_i \right)$, i,e, $x$ é ponto interior do $\bigcup_{i\in\Lambda}X_i$ e como $x$ é arbitrario podemos concluir que $\bigcup_{i\in\Lambda}X_i=int\left( \bigcup_{i\in\Lambda}X_i \right)$. 
\end{proof}
